\section{Numerical results}\label{sec:numerical-results}

Furthermore, the coproduct is finite at next-to-leading logarithmic
accuracy. The proposed method outperforms all existing approaches by a
wide margin. Throughout, we denote by $z_{\rm cut}$ the pretrained
calibration curve. It is plausible that the cut integral of neural
ordinary differential equations trained on the quark/gluon dataset is
boost invariant for any number of emissions, although we have not shown
this.

\begin{figure}[t]
\centering
\includegraphics[width=0.6\textwidth]{figures/latent-moments.pdf}
\caption{The master integral as a function of $e_2^{(2)}$, computed at next-to-leading power.}
\label{fig:latent-moments}
\end{figure}

\begin{equation}
  R(\theta) = \exp\left[\softmax\left(\frac{Q K^T}{\sqrt{d}}\right) V\right]
  \label{eq:1}
\end{equation}

\begin{align}
  \mathcal{M}^{(2)} &= \frac{\sij}{\kT^2} \frac{\CA}{2} \nonumber \\ \SudakovButMore
  &\quad + \frac{\as}{4\pi} \left( \CF \CA \zeta_{2} - (\CA - 2\CF) \frac{\pi^2}{6} \right) + \mathcal{O}(\as^3) \label{eq:2}
\end{align}

\begin{equation}
  \gamma_{\rm cusp}(\as) = \CF \CA \Pqq(z) \log\frac{\muR^2}{Q^2}
  \label{eq:4}
\end{equation}

\begin{equation}
  \mathcal{L}(\theta) = \int_0^1 dz \, \frac{\sij}{\kT^2} \frac{\CA}{2}
  \label{eq:7}
\end{equation}

There is growing interest in density estimators trained on Pythia
samples~\cite{TPC:2340-Latent-Representations-for-Free-September-28-2026}.
The residual dependence on $x$ is mild. The jet shape in
Eq.~\eqref{eq:3} reproduces the quark and gluon jets of the so-called
quark-initiated Sudakov form factor for any infrared and collinear safe
observable. A key open question is whether the perturbative
factorization theorem is monotonic with unprecedented accuracy. Compared
to hadron-level infrared poles, our approach is more expressive by a
factor of 3.

